\section{Estimation}\label{sec:estimation}

% ---------------------------------------------------------------------------
\subsection{Marginal likelihood via Laplace approximation}\label{sec:est_laplace}

The marginal log-likelihood \eqref{eq:margloglik} integrates the latent field $\boldsymbol{\eta} \in \mathbb{R}^n$ out of the joint density $p(\by,\boldsymbol{\eta}\mid\btheta)$.
Because the ZINB observation model is nonlinear in $\boldsymbol{\eta}$, the integral lacks a closed form.
We approximate it using the Laplace method, which is the standard strategy for latent Gaussian models in spatial statistics \citep{rue_martino_chopin2009} and mixed-effects modelling \citep{brooks2017}, and is implemented efficiently in the \texttt{TMB} framework for \texttt{R} \citep{kristensen2016}.

Write the joint negative log-likelihood as
\begin{equation}\label{eq:neg_joint}
  -\log p(\by,\boldsymbol{\eta}\mid\btheta) = f(\boldsymbol{\eta};\,\btheta) \equiv
  -\sum_{i=1}^{n}\log f(y_i\mid\eta_i,\pi_i;\,\kappa)
  -\log p(\boldsymbol{\eta}\mid\bbeta,\rho,\sigma^2).
\end{equation}
The Laplace approximation to the marginal likelihood replaces the integral by its second-order expansion around the mode:
\begin{equation}\label{eq:laplace}
  \ell_{\mathrm{LA}}(\btheta) = -f(\hat{\boldsymbol{\eta}};\,\btheta) + \tfrac{n}{2}\log(2\pi) - \tfrac{1}{2}\log\det\!\bigl(\mathbf{H}(\hat{\boldsymbol{\eta}};\,\btheta)\bigr),
\end{equation}
where $\hat{\boldsymbol{\eta}}(\btheta) = \arg\min_{\boldsymbol{\eta}} f(\boldsymbol{\eta};\,\btheta)$ is the conditional mode and $\mathbf{H} = \nabla^2_{\boldsymbol{\eta}} f\,\big|_{\hat{\boldsymbol{\eta}}}$ is the Hessian of $f$ with respect to $\boldsymbol{\eta}$.

The inner optimisation (finding $\hat{\boldsymbol{\eta}}$) is a weighted least-squares problem in the latent field and is solved by Newton--Raphson iteration.
The outer optimisation (maximising $\ell_{\mathrm{LA}}$ over $\btheta$) is carried out by a box-constrained quasi-Newton method (\texttt{nlminb}), with gradients computed by automatic differentiation (AD) through both the inner and outer problems.

\paragraph{Exploiting sparsity.}
The Hessian $\mathbf{H}$ inherits the sparsity pattern of $\bW$: the SAR prior contributes $\sigma^{-2}(\mathbf{I}_n - \rho\bW)'(\mathbf{I}_n - \rho\bW)$, which is sparse, and the observation model contributes a diagonal term $\diag(d_1,\ldots,d_n)$ (since the data are conditionally independent given $\boldsymbol{\eta}$).
Hence $\mathbf{H}$ is sparse and its Cholesky factorisation can be computed in $O(n)$ operations for banded or nested-dissection orderings.
The log-determinant in \eqref{eq:laplace} is obtained as twice the sum of the log-diagonal of the Cholesky factor, avoiding any dense matrix inversion.
This makes the estimator scalable to $n = O(10^4)$ or beyond.

\paragraph{Log-determinant of the SAR filter.}
The prior density \eqref{eq:sar_prior} involves $|\det(\mathbf{I}_n - \rho\bW)|$.
Pre-computing the eigenvalues $\lambda_1,\ldots,\lambda_n$ of $\bW$ (a one-time $O(n^3)$ cost, or $O(n)$ for special structures), the log-determinant is
\begin{equation}\label{eq:logdet}
  \log|\det(\mathbf{I}_n - \rho\bW)| = \sum_{i=1}^{n}\log|1 - \rho\lambda_i|.
\end{equation}
For row-standardised $\bW$, $\lambda_1 = 1$ and $\lambda_n \geq -1$, so $\rho \in (-1,1)$ ensures non-singularity.
We enforce this via the reparametrisation $\rho = \tanh(\tilde{\rho})$.
When $n$ is large the dense eigendecomposition becomes the bottleneck (it requires $O(n^2)$ memory as well as $O(n^3)$ time), and we use instead
\begin{equation}\label{eq:logdet_sparse}
  \log|\det(\mathbf{I}_n - \rho\bW)| = \tfrac12\log\det\bigl[\mathbf{I}_n - \rho(\bW+\bW') + \rho^2\bW'\bW\bigr],
\end{equation}
whose argument, $(\mathbf{I}_n-\rho\bW)'(\mathbf{I}_n-\rho\bW)$, is sparse, symmetric and positive definite whenever $\mathbf{I}_n-\rho\bW$ is non-singular.
Its log-determinant is obtained from a sparse Cholesky factorisation at each evaluation of the objective and differentiated with respect to $\rho$ by \texttt{TMB}'s atomic sparse log-determinant, so that no $n\times n$ dense matrix is ever formed.
The two routes give identical estimates and standard errors; the sparse one is used in the application of Section~\ref{sec:application}, where $n = 32{,}844$.

% ---------------------------------------------------------------------------
\subsection{Estimation algorithm}\label{sec:est_algorithm}

Algorithm~\ref{alg:zisar} summarises the full procedure, implemented in \texttt{TMB} \citep{kristensen2016}.

\begin{algorithm}[htbp]
\caption{Marginal-likelihood estimation of the ZISAR model via Laplace approximation.}\label{alg:zisar}
\begin{algorithmic}[1]
\Require Data $(\by, \bX, \bZ, \bW)$; eigenvalues $\lambda_1,\ldots,\lambda_n$ of $\bW$; convergence tolerance $\epsilon$.
\Ensure Parameter estimates $\hat{\btheta}$, standard errors, marginal log-likelihood.
\State \textbf{Pre-compute:} Eigenvalues of $\bW$ (for large $n$, the sparse matrices $\bW+\bW'$ and $\bW'\bW$ of \eqref{eq:logdet_sparse} instead); sparse Cholesky ordering for $(\mathbf{I}-\rho\bW)'(\mathbf{I}-\rho\bW)$.
\State \textbf{Initialise:} $\bbeta^{(0)}$ from an aspatial Poisson GLM; $\bgamma^{(0)} = (\logit(\bar{p}_0), 0, \ldots, 0)'$ with $\bar{p}_0$ the observed share of zeros; $\rho^{(0)}=0$, $\kappa^{(0)}=1$, $\sigma^{(0)}=0.5$; $\boldsymbol{\eta}^{(0)} = \bX\bbeta^{(0)}$.
\State \textbf{Construct TMB objective:} Define $f(\boldsymbol{\eta};\,\btheta)$ as the negative joint log-density \eqref{eq:neg_joint}, with $\boldsymbol{\eta}$ declared as random effects.
\Repeat
  \State \textbf{Inner loop:} Find $\hat{\boldsymbol{\eta}}(\btheta^{(s)}) = \arg\min_{\boldsymbol{\eta}} f(\boldsymbol{\eta};\,\btheta^{(s)})$ via Newton--Raphson with sparse Cholesky.
  \State \textbf{Laplace:} Evaluate $\ell_{\mathrm{LA}}(\btheta^{(s)})$ and $\nabla_{\btheta}\ell_{\mathrm{LA}}(\btheta^{(s)})$ by automatic differentiation.
  \State \textbf{Outer step:} Update $\btheta^{(s+1)}$ by a box-constrained quasi-Newton step (PORT routine \texttt{nlminb}).
\Until{$\|\btheta^{(s+1)} - \btheta^{(s)}\| < \epsilon$}
\State $\hat{\btheta} \leftarrow \btheta^{(s+1)}$, $\hat{\boldsymbol{\eta}} \leftarrow \hat{\boldsymbol{\eta}}(\hat{\btheta})$.
\State Compute the Hessian of $-\ell_{\mathrm{LA}}$ w.r.t.\ $\btheta$ by differentiating the AD gradient; invert for the variance--covariance matrix $\hat{\mathbf{V}}$.
\State \Return $\hat{\btheta}$, $\sqrt{\diag(\hat{\mathbf{V}})}$, $\ell_{\mathrm{LA}}(\hat{\btheta})$.
\end{algorithmic}
\end{algorithm}

\paragraph{Practical details.}
The reparametrisation $\rho = \tanh(\tilde{\rho})$, $\kappa = \exp(\tilde{\kappa})$, and $\sigma = \exp(\tilde{\sigma})$ maps the constrained parameter space to $\mathbb{R}^{\dim(\btheta)}$. For numerical stability the optimiser nonetheless works within wide box constraints, $\tilde\rho\in[-4,4]$ ($|\rho|\le0.9993$), $\kappa\in[0.01,500]$ and $\sigma\in[0.001,10]$, and is restarted from up to three perturbed starting points if it fails to converge; fits ending at the upper bound for $\kappa$ are flagged (Section~\ref{sec:montecarlo}).
Standard errors on the original scale are obtained via the delta method (handled automatically by \texttt{TMB}'s \texttt{sdreport()}).

% ---------------------------------------------------------------------------
\subsection{Asymptotic properties}\label{sec:est_asymptotics}

\begin{proposition}[Consistency]\label{prop:consistency}
Under Assumptions~\ref{ass:W}--\ref{ass:identification}, the Laplace-approximate MLE $\hat{\btheta}_n$ satisfies $\hat{\btheta}_n \plim \btheta_0$ as $n \to \infty$.
\end{proposition}

\begin{proof}[Proof sketch]
The Laplace approximation error is $O(n^{-1})$ \citep{tierney_kadane1986}, so $n^{-1}\ell_{\mathrm{LA}}(\btheta) - n^{-1}\ell(\btheta) \to 0$ uniformly over $\Theta$.
The true marginal log-likelihood $n^{-1}\ell(\btheta)$ converges uniformly to $\E[\ell_1(\btheta)]$, which is uniquely maximised at $\btheta_0$ (Assumption~\ref{ass:identification}).
Consistency follows from the standard extremum-estimator argument \citep{newey_mcfadden1994}.
A complete proof is in Appendix~\ref{app:proofs}.
\end{proof}

\begin{proposition}[Asymptotic normality]\label{prop:normality}
Under Assumptions~\ref{ass:W}--\ref{ass:identification} and standard smoothness conditions on $\ell_{\mathrm{LA}}$,
\[
  \sqrt{n}\,(\hat{\btheta}_n - \btheta_0) \dlim N\!\bigl(\mathbf{0},\,\mathcal{I}(\btheta_0)^{-1}\bigr),
\]
where $\mathcal{I}(\btheta_0) = -\lim_{n\to\infty} n^{-1}\nabla^2_{\btheta}\ell(\btheta_0)$ is the Fisher information per observation.
\end{proposition}

\begin{proof}[Proof sketch]
The result follows from the mean-value expansion of the Laplace-approximate score, the CLT for spatially dependent arrays under the bounded-dependence condition (Assumption~\ref{ass:W}), and the vanishing Laplace-approximation error.
Details are in Appendix~\ref{app:proofs}.
\end{proof}

% ---------------------------------------------------------------------------
\subsection{Direct, indirect, and total effects}\label{sec:est_effects}

In the ZISAR model, regression coefficients $\bbeta$ do not directly represent marginal effects because of the interplay between the exponential link, the spatial multiplier $\mathbf{S}(\rho)$, and the zero-inflation mixture.
Following \citet{lesage_pace2009} and adapting to the nonlinear ZI structure, we derive the effects as follows.

At the estimated mode $\hat{\boldsymbol{\eta}}$, the fitted count mean and inflation probability are
\begin{equation}\label{eq:fitted}
  \hat{\mu}_i = \exp(\hat{\eta}_i), \qquad \hat{\pi}_i = \logit^{-1}(\mathbf{z}_i'\hat{\bgamma}),
\end{equation}
and the marginal expectation is $\E[y_i] \approx (1 - \hat{\pi}_i)\hat{\mu}_i$.

\paragraph{Effects through the count process.}
Consider the $r$-th column of $\bX$.
A change in $x_{jr}$ (unit $j$, covariate $r$) propagates to $\eta_i$ for all $i$ through the spatial multiplier:
\begin{equation}\label{eq:deta_dx}
  \frac{\partial\boldsymbol{\eta}}{\partial x_{\cdot r}} = \mathbf{S}(\rho)\,\beta_r.
\end{equation}
The resulting effect on $\E[\by]$ is, stacking all $(i,j)$ pairs into an $n \times n$ matrix,
\begin{equation}\label{eq:effects_count}
  \mathbf{M}_r = \diag\!\bigl((\mathbf{1}-\hat{\bpi})\odot\hat{\bmu}\bigr)\;\mathbf{S}(\rho)\;\beta_r,
\end{equation}
where $\odot$ denotes element-wise multiplication.

\paragraph{Effects through the inflation process.}
If the $r$-th covariate also enters $\bZ$, the inflation probability responds:
\begin{equation}\label{eq:effects_zero}
  \frac{\partial\pi_i}{\partial z_{ir}} = \pi_i(1-\pi_i)\,\gamma_r,
\end{equation}
contributing an additional channel
\begin{equation}\label{eq:effects_zero_channel}
  \mathbf{M}_r^{(\pi)} = -\diag(\hat{\bmu})\;\diag\!\bigl(\hat{\bpi}\odot(\mathbf{1}-\hat{\bpi})\bigr)\;\gamma_r\;\mathbf{I}_n.
\end{equation}
Under the restricted ZISAR, $\mathbf{g} = \mathbf{0}$ implies no spatial feedback in the inflation process, so $\mathbf{M}_r^{(\pi)}$ is diagonal---the zero-inflation channel produces only direct effects.

\paragraph{Summary measures.}
The average direct, indirect, and total effects are defined as:
\begin{align}
  \overline{M}_r^{\,\mathrm{direct}} &= n^{-1}\,\tr\!\bigl(\mathbf{M}_r + \mathbf{M}_r^{(\pi)}\bigr), \label{eq:avg_direct} \\
  \overline{M}_r^{\,\mathrm{total}}  &= n^{-1}\,\mathbf{1}_n'\,\bigl(\mathbf{M}_r + \mathbf{M}_r^{(\pi)}\bigr)\,\mathbf{1}_n, \label{eq:avg_total} \\
  \overline{M}_r^{\,\mathrm{indirect}} &= \overline{M}_r^{\,\mathrm{total}} - \overline{M}_r^{\,\mathrm{direct}}. \label{eq:avg_indirect}
\end{align}
The indirect effects arise entirely through the count-process channel (the spatial multiplier $\mathbf{S}(\rho)$), since the inflation channel is diagonal under the restricted specification.

\begin{proposition}[Effects decomposition]\label{prop:effects}
Under the restricted ZISAR ($\mathbf{g}=\mathbf{0}$), the $n \times n$ matrix of effects of the $r$-th count covariate on $\E[\by]$ is given by \eqref{eq:effects_count}.
These effects are:
\begin{enumerate}
\item[(i)] strictly positive (negative) if $\beta_r > 0$ ($\beta_r < 0$) and $\rho \in [0,1)$;
\item[(ii)] the indirect effect vanishes if and only if $\rho = 0$;
\item[(iii)] monotonically increasing in $|\rho|$ for fixed $\bbeta$.
\end{enumerate}
\end{proposition}

\begin{proof}
Follows from the non-negativity of all elements of $\mathbf{S}(\rho) = \sum_{s=0}^{\infty}\rho^s\bW^s$ when $\rho \geq 0$ and $\bW \geq 0$, and from the strict positivity of $(1-\hat{\pi}_i)\hat{\mu}_i > 0$.
\end{proof}

Standard errors for the average effects are obtained by simulation from the asymptotic distribution: draw $B$ replications $\btheta^{(b)} \sim N(\hat{\btheta}, \hat{\mathbf{V}})$, compute $\overline{M}_r^{(b)}$ at each draw, and report the standard deviation and percentile-based confidence intervals.
